\renewcommand{\abstractname}{\large\bfseries\color{primary} Resumo}
\begin{abstract}
\noindent As demandas severas de componentes de reentrada atmosférica e veículos hipersónicos exigem novos materiais capazes de manter integridade estrutural sob elevadas solicitações aerotérmicas, oxidativas e mecânicas. As Ligas de Alta Entropia (HEAs) e, em particular, as Ligas de Alta Entropia Refratárias (RHEAs) surgem como uma das fronteiras mais promissoras da engenharia de materiais de ultra-alta temperatura (UHT). Este trabalho apresenta um estudo cienciométrico e de revisão crítica estruturado sobre um modelo metodológico replicável. A partir de uma estratégia de busca combinatória automatizada via script Python cruzando 42 descritores de ligas e 22 descritores de severidade aerotérmica nos campos \texttt{TI OR AB OR AK} na \emph{Web of Science Core Collection} (período 2004--2026), consolidou-se um corpus macroscópico de 192 documentos de alta especialização. O estudo implementa rotinas de higienização de dados (\emph{thesaurus}) e mitigação do efeito de homonímia na desambiguação autoral (AND). Concomitantemente à caracterização quantitativa das frentes de publicação via Bibliometrix e VOSviewer, realiza-se a síntese qualitativa dos mecanismos de endurecimento, estabilidade de fases BCC/B2, limites de escoamento até 2000 \textdegree{}C e o compromisso crítico entre ductilidade a frio e resistência à oxidação/\emph{pesting} em regimes de Mach 5+.

\vspace{0.8em}
\noindent {\color{primary}\textbf{Palavras-chave:}} Ligas de Alta Entropia; Ligas de Alta Entropia Refratárias (RHEAs); Voo Hipersónico; Cienciometria; Arquitetura de Duplo Corpus (BACR); Desambiguação Autoral; Bibliometria.
\end{abstract}
